\RequirePackage[cmintegrals,cmbraces,ebgaramond]{newtxmath}
\InputIfFileExists{../global.src}\relax\relax

\iffull
\title[Fünftes Tutorium -- Übungsblatt 5]{Raute-Hashtag-Matrizen\\\small Mathemann 5}
\date{\sffamily KW 48}

\usepackage{tikzpingus}

\usetikzlibrary{decorations.pathreplacing,decorations.text}

\hypersetup{colorlinks=false}

% \lstset{add to literate={Vector}{{{\solGet{keywordC}{Vector}}}}6}

\begin{document}
\Titlepage{5}
\fi
\section{Präsenzaufgabe}
\begin{frame}[fragile, c]{Präsenzaufgabe, zum Auflockern}
    \begin{aufgabe}{Flippn' Amazing -- Watch me Flip-Flip\ldots}
        \vspace*{-.125\topsep-\baselineskip}\par\small\pause Implementieren Sie die Methoden \bjava{public static void flipInPlace(char[] x)} und \bjava{public static char[] flipInCopy(char[] x)}, welche beide die Groß- und Kleinschreibung in den Arrays umkehren (nur ASCII).\pause\ \T{flipInPlace} soll das übergeben Array modifizieren, \T{flipInCopy} soll \T{x} aber nicht verändern und eine \say{geflippte} Kopie zurückliefern.

        \pause Zum Erstellen einer Kopie können Sie \bjava{char[] copy = x.clone()} verwenden.\pause\ Nutzen Sie aber nicht \bjava{Character::toUpperCase} oder vergleichbare Methoden,\pause\ sondern machen Sie sich die Unicode-Kodierung von \bjava{char} zunutze (\(A\)--\(Z\) ist \(65\)--\(90\), \(a\)--\(z\) ist \(97\)--\(122\)).

        \pause Beantworten Sie weiter:\vspace*{-.66\topsep} \begin{enumerate}[<+(1)->]
            \itemsep0pt
            \item Was sind die Vor- und Nachteile beider Ansätze?
            \item Warum reicht keine einfache Zuweisung zur Kopie des Arrays?
            \item Warum funktioniert kein for-each, wenn man die Einträge modifizieren möchte?
        \end{enumerate}
        \onslide<1->
    \end{aufgabe}
\end{frame}

\begin{frame}[t,fragile]{Präsenzaufgabe - Lösung}
    \begin{itemize}[<+(1)->]
        \item We do se flippin place (mit \T{x} als \T{flipEm}):
\begin{plainjava}
!*\onslide<3->*!public static void flipInPlace(char[] flipEm) {
!*\onslide<4->*!  int shift = 'a' - 'A';
!*\onslide<5->*!  for (int i = 0; i < flipEm.length; i++) {
!*\onslide<6->*!    if ('A' <= flipEm[i] && flipEm[i] <= 'Z') {
!*\onslide<7->*!      flipEm[i] += shift;
!*\onslide<6->*!    } else if ('a' <= flipEm[i] && flipEm[i] <= 'z') {
!*\onslide<8->*!      flipEm[i] -= shift;
!*\onslide<6->*!    }
!*\onslide<5->*!  }
!*\onslide<3->*!}
\end{plainjava}
    \end{itemize}
\end{frame}

\begin{frame}[t,fragile]{Präsenzaufgabe - Lösung}
    \begin{itemize}[<+(1)->]
        \item Und nun mit der Kopie! Wir könnten das ja nochmal schreiben\ldots
\begin{plainjava}
!*\onslide<3->*!public static char[] flipInCopy(char[] flipEm) {
!*\onslide<4->*!    char[] copy = flipEm.clone();
!*\onslide<5->*!    flipInPlace(copy);
!*\onslide<6->*!    return copy;
!*\onslide<3->*!}
\end{plainjava}
    \end{itemize}
\end{frame}

\begin{frame}[t,fragile]{Präsenzaufgabe - Lösung}
    \begin{itemize}[<+(1)->]
        \itemsep12pt
        \item Vorteile einer Kopie (gegenüber in-place): \begin{itemize}
            \item Nebeneffekte von Methoden sind generell schlecht! \say{out-Parameter} sollten vermieden werden (wo nicht unbedingt notwendig).
        \end{itemize}
        \item Nachteile einer Kopie (gegenüber in-place): \begin{itemize}
            \item Der benötigte Speicher wird verdoppelt (was, wenn das Array riesig ist?)
        \end{itemize}
        \item Was ist nun besser?
            \begin{itemize}
                \item In der Regel die Variante mit Kopie. Die unerwarteten/ungewollten Seiteneffekte können Kaskaden schwer zu findender Fehler verursachen.
                \item Sollten wirklich große Arrays erwartet werden, sollte man sich unter Umständen ein grundlegend anderes System überlegen.
            \end{itemize}
    \end{itemize}
\end{frame}

\begin{frame}[t,fragile]{Präsenzaufgabe - Lösung}
    \begin{itemize}[<+(1)->]
        \itemsep8pt
        \item Warum reicht \bjava{char[] clone = x;} nicht für eine Kopie?\smallskip\\*
            \pause Arrays sind komplexe Datenstrukturen (\say{ReferenceType}).\pause\ Mit \bjava{clone = x} lassen wir \T{clone} auf die gleiche Speicheradresse wie \T{x} zeigen.\pause\ Wir verändern dasselbe Array.
        \item Warum keine for-each Schleife?\smallskip\\*
            \pause Damit erhalten wir nur die Werte an der jeweiligen Stelle, aber nicht den Index den es zu verändern gilt. Der Iterator erzeugt eine Kopie:\pause
\begin{plainjava}
for(char c : flipEm) { /* hier haben wir nur c */ }
\end{plainjava}
        \pause Wir könnten zwar einen Zeiger inkrementieren, aber wer garantiert uns dann die Reihenfolge der Iteration?\pause\ Und selbst wenn das passt\ldots\pause\ Im Endeffekt haben wir dann wieder eine for-Schleife.
    \end{itemize}
\end{frame}


\section{Übungsblatt 5}
\subsection{Aufgabe 1}
\begin{frame}[t,fragile]{Übungsblatt 5 - Aufgabe 1 a)}
    \begin{itemize}[<+(1)->]
        \item Der volle Code befindet sich hier: \only<2->{\textattachfile{\curpath DotProduct.java}{DotProduct.java}}.
        \item Wir berechnen \(d = a \cdot b = (\begin{smallmatrix}
            a_1 & \ldots & a_n
        \end{smallmatrix}) \cdot (\begin{smallmatrix}
            b_1 & \ldots & b_n
        \end{smallmatrix}) = \sum_{i = 1}^n a_i \cdot b_i\)
        \item Oder als Java-Code:
\begin{plainjava}
!*\onslide<5->*!int d = 0;
!*\onslide<6->*!for(int i = 0; i < length; i++) {
!*\onslide<6->*!    d += vector1[i] * vector2[i];
!*\onslide<6->*!}
!*\onslide<5->*!
!*\onslide<7->*!System.out.println("Vector1: " + vectorToString(vector1));
!*\onslide<7->*!System.out.println("Vector2: " + vectorToString(vector2));
!*\onslide<5->*!
!*\onslide<8->*!System.out.println("Das Skalarprodukt ist: " + d);
\end{plainjava}
    \end{itemize}
\end{frame}

\begin{frame}[t,fragile]{Übungsblatt 5 - Aufgabe 1 a)}
    \begin{itemize}[<+(1)->]
        \item Nun noch Ausgabe der Vektoren:
\begin{plainjava}
!*\onslide<3->*!public static String vectorToString(int[] vector) {
!*\onslide<4->*!    String result = "(";
!*\onslide<5->*!    for (int i = 0; i < vector.length; i++) {
!*\onslide<6->*!        if (i > 0) result += ", ";
!*\onslide<7->*!        result += vector[i];
!*\onslide<5->*!    }
!*\onslide<8->*!    return result + ")";
!*\onslide<3->*!}
\end{plainjava}
    \end{itemize}
\end{frame}

\iffull
\rExecute{javac DotProduct.java}
\begin{frame}[t,fragile]{Übungsblatt 5 - Aufgabe 1 a)}
    \begin{itemize}[<+(1)->]
        \item Als beispielhafte Ausgabe für \rbash[:\pause]{java DotProduct}
    \end{itemize}
\end{frame}
\fi

\begin{frame}[t,fragile]{Übungsblatt 5 - Aufgabe 1 b)}
    \begin{itemize}[<+(1)->]
        \item Der volle Code befindet sich hier: \only<2->{\textattachfile{\curpath Transpose.java}{Transpose.java}}.
        \item Wir vertauschen! Aus \(a_{i, j}\) wird \(a'_{j, i}\).
        \item Oder als Java-Code:
\begin{plainjava}
!*\onslide<5->*!int[][] transposed = new int[M][N];
!*\onslide<6->*!for (int i = 0; i < M; i++) {
!*\onslide<7->*!    for (int j = 0; j < N; j++) {
!*\onslide<8->*!        transposed[i][j] = matrix[j][i];
!*\onslide<8->*!        System.out.print(transposed[i][j] + " ");
!*\onslide<7->*!    }
!*\onslide<9->*!    System.out.println("");
!*\onslide<6->*!}
\end{plainjava}
    \end{itemize}
\end{frame}

\iffull
\rExecute{javac Transpose.java}
\begin{frame}[t,fragile]{Übungsblatt 5 - Aufgabe 1 b)}
    \def\rbashcurrentinject{multicols=2}%
    \begin{itemize}[<+(1)->]
        \item \lstfs{9}Als beispielhafte Ausgabe für \rbash[:\pause]{java Transpose}
    \end{itemize}
\end{frame}
\fi

\subsection{Aufgabe 2}
\begin{frame}[t,fragile]{Übungsblatt 5 - Aufgabe 2}
    \begin{itemize}[<+(1)->]
        \item Der volle Code befindet sich hier: \only<2->{\textattachfile{\curpath MatrixMultiplication.java}{MatrixMultiplication.java}}.
        \item Zunächst eine zweite Matrix:
\begin{plainjava}
!*\onslide<4->*!int[][] matrix2 = new int[M][N];
!*\onslide<5->*!for(int i = 0; i < M; i++) {
!*\onslide<6->*!    for (int j = 0; j < N; j++) {
!*\onslide<7->*!        matrix2[i][j] = (int) (Math.random() * 10);
!*\onslide<8->*!        System.out.print(matrix2[i][j] + " ");
!*\onslide<6->*!    }
!*\onslide<9->*!    System.out.println("");
!*\onslide<5->*!}
!*\onslide<10->*!System.out.println("");
\end{plainjava}
    \end{itemize}
\end{frame}

\begin{frame}[t,fragile]{Übungsblatt 5 - Aufgabe 2}
    \begin{itemize}[<+(1)->]
        \item Und nun multiplizieren wir! Da wir davon ausgehen können, dass es klappt, gilt:
        \onslide<3->{\(c_{i, j} = \sum_{k = 1}^m a_{i, k} \cdot b_{k, j}\).}\lstfs{10}
\begin{plainjava}
!*\onslide<4->*!int[][] result = new int[N][N];
!*\onslide<5->*!for(int i = 0; i < N; i++) {
!*\onslide<6->*!    for (int j = 0; j < N; j++) {
!*\onslide<7->*!        int tmp = 0;
!*\onslide<8->*!        for (int k = 0; k < M; k++) {
!*\onslide<9->*!            tmp += matrix1[i][k] * matrix2[k][j];
!*\onslide<8->*!        }
!*\onslide<10->*!        result[i][j] = tmp;
!*\onslide<11->*!        System.out.print(result[i][j] + " ");
!*\onslide<6->*!    }
!*\onslide<12->*!    System.out.println();
!*\onslide<5->*!}
\end{plainjava}
    \end{itemize}
\end{frame}

\iffull
\rExecute{javac MatrixMultiplication.java}
\begin{frame}[t,fragile]{Übungsblatt 5 - Aufgabe 2}
    \def\rbashcurrentinject{multicols=2}%
    \begin{itemize}[<+(1)->]
        \item \lstfs{8}Als beispielhafte (formatierte) Ausgabe für \rbash[:\pause]{java MatrixMultiplication}
    \end{itemize}
\end{frame}
\fi

\subsection{Zusatzaufgabe 2/6}
\begin{frame}[fragile]{Übungsblatt 5 - Zusatzaufgabe a)}
    \begin{itemize}[<+(1)->]
        \item Wir machen eine Raute (\only<2->{\textattachfile{\curpath Rhombus.java}{Rhombus.java}}).
        \item Naiv, können wir die bekannte Dreiecksstruktur nehmen und ein Dreieck für die obere und eines für die untere Hälfte basteln:
\begin{plainjava}
!*\onslide<4->*!// oben
!*\onslide<4->*!for(int i = 1; i <= N; i++) {
!*\onslide<5->*!    for(int j = 1; j <= N-i; j++) System.out.print(" ");
!*\onslide<4->*!
!*\onslide<6->*!    for(int j = 1; j <= i; j++) {
!*\onslide<7->*!        if((j == 1) || (j == i)) System.out.print("* ");
!*\onslide<8->*!        else System.out.print("  ");
!*\onslide<6->*!    }
!*\onslide<9->*!    System.out.println("");
!*\onslide<4->*!}
\end{plainjava}
    \end{itemize}
\end{frame}

\begin{frame}[fragile]{Übungsblatt 5 - Zusatzaufgabe a)}
    \begin{itemize}[<+(1)->]
        \item Für unten spiegeln wir:
\begin{plainjava}
!*\onslide<3->*!for(int i = N; i > 1; i--) {
!*\onslide<4->*!    for(int j = i-1; j < N; j++) System.out.print(" ");
!*\onslide<3->*!
!*\onslide<5->*!    for(int j = i; j >= 2; j--) {
!*\onslide<6->*!        if((j == 2) || (j == i)) System.out.print("* ");
!*\onslide<7->*!        else System.out.print("  ");
!*\onslide<5->*!    }
!*\onslide<8->*!    System.out.println("");
!*\onslide<3->*!}
\end{plainjava}
    \end{itemize}
\end{frame}

\begin{frame}[fragile]{Übungsblatt 5 - Zusatzaufgabe a)}
    \begin{itemize}[<+(1)->]
        \item Negative Zahlen erlauben es, beide Hälften in einer Schleife zu behandeln:
\begin{plainjava}
!*\onslide<3->*!for(int i = N-1; i > -N; i--) {
!*\onslide<4->*!    for(int j = 1; j <= Math.abs(i); j++)
!*\onslide<5->*!        System.out.print(" ");
!*\onslide<3->*!
!*\onslide<6->*!    for(int j = 1; j <= N - Math.abs(i); j++) {
!*\onslide<7->*!        if((j == 1) || (j == N - Math.abs(i)))
!*\onslide<8->*!            System.out.print("* ");
!*\onslide<7->*!        else
!*\onslide<8->*!            System.out.print("  ");
!*\onslide<6->*!    }
!*\onslide<9->*!    System.out.println("");
!*\onslide<3->*!}
\end{plainjava}
    \end{itemize}
\onslide<10->{\begin{tikzpicture}[overlay,remember picture]
    \node[above left=.55cm,xshift=.25cm] at(current page.south east) {\only<10->{\textattachfile{\curpath RhombusNegative.java}{RhombusNegative.java}}};
\end{tikzpicture}}
\end{frame}

\iffull
\rExecute{javac RhombusNegative.java}
\begin{frame}[t,fragile]{Übungsblatt 5 - Zusatzaufgabe a)}
    \lstfs{9}\begin{itemize}[<+(1)->]
        \item So ergibt sich beispielsweise \rbash[:\pause]{java RhombusNegative 9}
    \end{itemize}
\end{frame}

\begin{frame}[c,plain]{}
\tikzpicture[overlay,remember picture]
\fill[pingu@black] (current page.south east) rectangle (current page.north west);
\onslide<2->{\node[pingu@white,above,scale=1.66] (c) at (current page.center) {So tell me\ldots};}
\onslide<3->{\node[pingu@white,below,scale=1.66] at (c.south) {What is it, you \textit{desire}?};}
\onslide<4->{
    \path (current page.center)++(50:3cm) node[pingu@white,below,scale=1.3,rotate=-20] {\bfseries Math!};
}
\onslide<5->{
    \path (current page.center)++(-210:4.25cm) node[pingu@white,below,scale=1.9,rotate=15] {\bfseries Math!};
}
\onslide<6->{
    \path (current page.center)++(-60:3.25cm) node[pingu@white,below,scale=1.45,rotate=45] {\bfseries Math!};
}
\onslide<7->{
    \path (current page.center)++(215:3.75cm) node[pingu@white,below,scale=1.85,rotate=-5] {\bfseries Math!};
}
\onslide<8->{
    \path (current page.center)++(24:5.95cm) node[pingu@white,below,scale=0.75,rotate=-19] {\bfseries \LaTeX-Penguins!};
}
\onslide<9->{
    \path (current page.center)++(-16:6.75cm) node[pingu@white,below,scale=1.65,rotate=-15] {\bfseries Math!};
}
\endtikzpicture
\end{frame}

\newsavebox\whitepingu
\savebox\whitepingu{\tikz\pingu[body=pingu@white,bill color=pingu@white,feet color=pingu@white,eyes color=pingu@white,heart=pingu@black];}
\begin{frame}[c,plain]{}
\tikzpicture[overlay,remember picture]
\fill[pingu@black] (current page.south east) rectangle (current page.north west);
\node[pingu@white,scale=1.66] (c) at (current page.center) {So Math it is\ldots};
\onslide<2->{\node[above left=3mm,scale=.75] at(current page.south east) {\usebox\whitepingu};}
\endtikzpicture
\end{frame}
\fi

\begin{frame}[c,plain]{}
\tikzpicture[overlay,remember picture]
\fill[pingu@black] (current page.south east) rectangle (current page.north west);
\only<1-2>{\node[below right=3mm,pingu@white] at(current page.north west) {\large\ttfamily \textcolor{purple!80!white}{Love}~\faAngleRight\ \textcolor{purple!40!white}{Real \scalebox{.8}{\faHeart}}~\faAngleRight\ \textcolor{green!40!white}{Math}};}
\node[above left=3mm,pingu@white,scale=.8] (pn) at(current page.south east) {\footnotesize\insertframenumber\thinspace/\thinspace\totalframes};
\only<1-2>{\node[above right=3mm,pingu@white,scale=.33] at(current page.south west) {\footnotesize Manuell Layouts im Foliensatz zu ändern macht Spaß. Wirkt aber fast, als hätte ich das hier geklaut :/ Das hab ich nich!};}

\onslide<2->{\scope[shift={([yshift=-6pt]current page.center)},pingu@white,scale=.675,every path/.append style={line cap=round}]

\draw[->] (-5.5,0) -- (5.5,0);
\draw[->] (0,-5.5) -- (0,5.5);
\only<2>{\foreach \i in {-5,-4,...,-1,1,2,...,5} {
    \draw (\i,.1) -- ++(0,-.2) node[below,scale=.875] {\(\i\)};
    \draw (.1,\i) -- ++(-.2,0) node[left,scale=.875] {\(\i\)};
}}
\only<3-|handout:0>{\foreach \i in {-5,-4,...,-1,1,2,...,5} {
    \draw (\i,.1) -- ++(0,-.2);
    \draw (.1,\i) -- ++(-.2,0);
}}

\only<4-5>{
    \colorlet{major}{paletteB!20!white}
    \only<handout:0>{\node[below left,outer sep=0pt,inner sep=0pt,major] at (5.5,5.5) {\vphantom{$\sqrt{y^2}$}\(x + y \leq 1\)};}
    \only<5>{\fill[fill opacity=.33,major,rounded corners=.5pt] (-4.5,5.5) -- (5.5,-4.5) -- (5.5,-5.5) -| (-5.5,5.5) -- cycle;
    \draw[major,thick] (-4.5,5.5) -- (5.5,-4.5);}
}

\only<6-7>{
    \colorlet{major}{paletteB!20!white}
    \only<handout:0>{\node[below left,outer sep=0pt,inner sep=0pt,major] at (5.5,5.5) {\vphantom{$\sqrt{y^2}$}\(\sqrt{x^2 + y^2} \leq 3\)};}
    \only<7>{\filldraw[major,fill opacity=.33,thick] (0,0) circle [radius=3cm];}
}

\only<8>{
    \colorlet{major}{paletteB!20!white}
    \only<handout:0>{\node[below left,outer sep=0pt,inner sep=0pt,major] at (5.5,5.5) {\vphantom{$\sqrt{y^2}$}\(\sqrt{x^2 + y^2} = 3\)};}
    \draw[major,thick] (0,0) circle [radius=3cm];
}

\only<9-10>{
    \colorlet{major}{paletteB!20!white}
    \only<handout:0>{\node[below left,outer sep=0pt,inner sep=0pt,major] at (5.5,5.5) {\vphantom{$\sqrt{y^2}$}\(\abs{x - y} + \abs{x + y} \leq 4\)};}
    \only<10>{\filldraw[major,fill opacity=.33,thick] (-2,-2) rectangle ++(4,4);}
}

\only<11-12>{s
    \colorlet{major}{paletteB!20!white}
    \only<handout:0>{\node[below left,outer sep=0pt,inner sep=0pt,major] at (5.5,5.5) {\vphantom{$\sqrt{y^2}$}\(\abs{x} + \abs{y} \leq 3\)};}
    \only<12>{\filldraw[major,fill opacity=.33,thick] (0,3) -- (-3,0) -- (0,-3) -- (3,0) -- cycle;}
}
\only<0|handout:1>{
\scope[every node/.append style={below left,scale=.8,outer xsep=0pt,inner xsep=0pt,major}]
    \node (1) at (5.5,5.5) {\(x + y \leq 1\)};
    \node (2) at (1.south east) {\(\sqrt{x^2 + y^2} \leq 3\)};
    \node (3) at (2.south east) {\(\abs{x - y} + \abs{x + y} \leq 4\)};
    \node (4) at (3.south east) {\(\abs{x} + \abs{x} \leq 4\)};
\scope[shift={([xshift=5mm,yshift=.9mm]1.east)},scale=.3]
    \draw[major] (-.75,1) -- (1,-.75);
    \fill[fill opacity=.33,major,rounded corners=.175pt] (-.75,1) -- (1,-.75) -- (1,-1) -| (-1,1) -- cycle;
\endscope
\scope[shift={([xshift=5mm,yshift=.5mm]2.east)},scale=.3]
    \filldraw[major,fill opacity=.33] (0,0) circle [radius=1cm];
\endscope
\scope[shift={([xshift=5mm,yshift=.9mm]3.east)},scale=.3]
    \filldraw[major,fill opacity=.33] (-.9,-.9) rectangle (.9,.9);
\endscope
\scope[shift={([xshift=5mm,yshift=.9mm]4.east)},scale=.3]
    \filldraw[major,fill opacity=.33] (0,.9) -- (-.9,0) -- (0,-.9) -- (.9,0) -- cycle;
\endscope
\endscope
}

\endscope}
\endtikzpicture
\end{frame}

\iffull
\begin{frame}[c,plain]{}
\tikzpicture[overlay,remember picture]
\fill[pingu@black] (current page.south east) rectangle (current page.north west);
\node[above left=3mm,pingu@white,scale=.8] (pn) at(current page.south east) {\footnotesize\insertframenumber\thinspace/\thinspace\totalframes};
\onslide<2->{\node[pingu@white,above=7.75mm,scale=1.66] (c) at (current page.center) {\textcolor{green!20!white}{Int}eressier\textcolor{teal!39!white}{t}?};}
\onslide<3->{\node[pingu@white,below,scale=1.66] at (c.south) {\textcolor{purple!30!white}{S}igned \textcolor{purple!30!white}{D}istance \textcolor{purple!30!white}{F}unctions};}

\onslide<4->{\node[pingu@white,below=2.25cm,scale=1] (h) at (c.south) {For the horde. For the rabbit hole!};}
\onslide<4->{\node[pingu@white,below,scale=1] at (h.south) {\disablehyper\textcolor{paletteB!30!white}{\href{https://iquilezles.org/}{Inigo Quilez}}~~\textbullet~~\textcolor{paletteB!30!white}{\href{https://www.youtube.com/watch?v=1iuLxJmQII0}{DeepSDF}}};}

\endtikzpicture
\end{frame}
\fi

\begin{frame}[fragile]{Übungsblatt 5 - Zusatzaufgabe a) \hfill Eleganter}
    \begin{itemize}
        \item<2-> Wir haben eine Menge zu der alle Punkte zählen für die gilt: \(\abs{x} + \abs{y} = N - 1\).
        \item<3-> \(N - 1\) rührt daher, da wir ja für \(N = 1\) sonst auch in \(y\)-Richtung je \(1\) hätten.\smallskip
\columns[c,onlytextwidth]
\column{.25cm}
\column{5.35cm}
\onslide<4->{\resizebox{5.1cm}!{\begin{tikzpicture}
    \draw[-Kite] (-4.5,0) -- (4.5,0);
    \draw[-Kite] (0,-4.5) -- (0,4.5);
    \foreach \i in {-4,-3,...,-1,1,2,...,4} {
        \draw (\i,.1) -- ++(0,-.2) node[below] {\(\i\)};
        \draw (.1,\i) -- ++(-.2,0) node[left] {\(\i\)};
    }
    \foreach \i/\col in {1/paletteA,2/paletteB,3/paletteC,4/paletteD} {
        \edef\nxt{\the\numexpr4+\i}
        \onslide<\nxt->{\draw[\col,thick] (0,\i) -- (-\i,0) -- (0,-\i) -- (\i,0) -- cycle;}
    }
    \only<8->{\node[above left,align=right,scale=1.33] at(current bounding box.south east) {\paletteA{\(N = 1\)}\\\paletteB{\(N = 2\)}\\\paletteC{\(N = 3\)}\\\paletteD{\(N = 4\)}};}
\end{tikzpicture}}}
\column{\linewidth-5.6cm}
\lstfs{10}\begin{plainjava}[lineskip=1.75pt]
!*\onslide<9->*!if(args.length != 1) System.exit(1);
!*\onslide<9->*!int N = Integer.parseInt(args[0]) - 1;
!*\onslide<10->*!// Diskretisieren
!*\onslide<10->*!for (int y = N; y >= -N; y--) {
!*\onslide<11->*!    for (int x = -N; x <= N; x++) {
!*\onslide<12->*!        if (Math.abs(y) + Math.abs(x) == N)
!*\onslide<12->*!            System.out.print("*");
!*\onslide<13->*!        else System.out.print(" ");
!*\onslide<11->*!    }
!*\onslide<14->*!    System.out.println();
!*\onslide<10->*!}
\end{plainjava}
\endcolumns
\onslide<15->{\begin{tikzpicture}[overlay,remember picture]
    \node[above left=.55cm,xshift=.25cm] at(current page.south east) {\only<15->{\textattachfile{\curpath RhombusElegant.java}{RhombusElegant.java}}};
\end{tikzpicture}}
    \end{itemize}
\end{frame}


\begin{frame}[fragile]{Übungsblatt 5 - Zusatzaufgabe b)}
    \begin{itemize}[<+(1)->]
        \item Die Musterlösung (\only<2->{\textattachfile{\curpath NestedRhombus.java}{NestedRhombus.java}}).\lstfs{7}
\begin{plainjava}
!*\onslide<3->*!for(int i = 1; i <= N; i++) {
!*\onslide<4->*!    for(int j = 1; j <= N-i; j++) System.out.print(" ");
!*\onslide<5->*!    if(i % 2 == 1) {
!*\onslide<6->*!        for (int j = 1; j <= i; j++) {
!*\onslide<7->*!            if (j % 2 == 1) System.out.print("* ");
!*\onslide<7->*!            else System.out.print("  ");
!*\onslide<6->*!        }
!*\onslide<5->*!    } else {
!*\onslide<8->*!        for (int j = 1; j <= i / 2; j++) { // Linke Hälfte
!*\onslide<9->*!            if (j % 2 == 1) System.out.print("* ");
!*\onslide<9->*!            else System.out.print("  ");
!*\onslide<8->*!        }
!*\onslide<10->*!        for (int j = i/2; j >= 1; j--) { // Rechte Hälfte
!*\onslide<11->*!            if (j % 2 == 1) System.out.print("* ");
!*\onslide<11->*!            else System.out.print("  ");
!*\onslide<10->*!        }
!*\onslide<5->*!    }
!*\onslide<12->*!    System.out.println("");
!*\onslide<3->*!}
\end{plainjava}
    \end{itemize}
\end{frame}

\begin{frame}[fragile]{Übungsblatt 5 - Zusatzaufgabe b)}
    \begin{itemize}[<+(1)->]
        \item Die Musterlösung der unteren Hälfte:\lstfs{7}
\begin{plainjava}
!*\onslide<3->*!for(int i = N-1; i >= 1; i--) {
!*\onslide<4->*!    for(int j = i; j < N; j++) System.out.print(" ");
!*\onslide<5->*!    if(i % 2 == 1) {
!*\onslide<6->*!        for (int j = 1; j <= i; j++) {
!*\onslide<7->*!            if (j % 2 == 1) System.out.print("* ");
!*\onslide<7->*!            else System.out.print("  ");
!*\onslide<6->*!        }
!*\onslide<5->*!    } else {
!*\onslide<8->*!        for (int j = 1; j <= i / 2; j++) { // Linke Hälfte
!*\onslide<9->*!            if (j % 2 == 1) System.out.print("* ");
!*\onslide<9->*!            else System.out.print("  ");
!*\onslide<8->*!        }
!*\onslide<10->*!        for (int j = i/2; j >= 1; j--) { // Rechte Hälfte
!*\onslide<11->*!            if (j % 2 == 1) System.out.print("* ");
!*\onslide<11->*!            else System.out.print("  ");
!*\onslide<10->*!        }
!*\onslide<5->*!    }
!*\onslide<12->*!    System.out.println("");
!*\onslide<3->*!}
\end{plainjava}
    \end{itemize}
\end{frame}

\begin{frame}[fragile]{Übungsblatt 5 - Zusatzaufgabe b) \hfill Eleganter}
    \begin{itemize}
        \item<2-> Wir möchten nun alle Rauten durch \(\abs{x} + \abs{y} \leq N - 1\).
        \item<3-> Da wir aber nur je die mit Längen \(- 4\) haben wollen: \(\abs{x} + \abs{y} \equiv N - 1 \mod 4\). \onslide<4->{Mathe ist Liebe.}
\columns[c,onlytextwidth]
\column{.25cm}
\column{5.35cm}
\onslide<5->{\resizebox{4.71cm}!{\begin{tikzpicture}
    \draw[-Kite] (-6.5,0) -- (6.5,0);
    \draw[-Kite] (0,-6.5) -- (0,6.5);
    \foreach \i in {-6,-5,...,-1,1,2,...,6} {
        \draw (\i,.1) -- ++(0,-.2) node[below] {\(\i\)};
        \draw (.1,\i) -- ++(-.2,0) node[left] {\(\i\)};
    }
    \foreach[count=\j] \i/\col in {2/paletteB,6/paletteB} {
        \edef\nxt{\the\numexpr5+\j}
        \onslide<\nxt->{\draw[\col,very thick] (0,\i) -- (-\i,0) -- (0,-\i) -- (\i,0) -- cycle;}
    }
    \only<7->{\node[above left,align=right,scale=1.33] at(current bounding box.south east) {\paletteB{\(N = 6\)}};}
\end{tikzpicture}}}
\column{\linewidth-5.6cm}
\lstfs{9}\begin{plainjava}[lineskip=1.75pt]
!*\onslide<8->*!if (args.length != 1) System.exit(1);
!*\onslide<8->*!int N = Integer.parseInt(args[0]) - 1;
!*\onslide<8->*!
!*\onslide<9->*!for (int y = N; y >= -N; y--) {
!*\onslide<10->*!    for (int x = -N; x <= N; x++) {
!*\onslide<11->*!       int radius = Math.abs(y) + Math.abs(x);
!*\onslide<12->*!       if (radius <= N && radius % 4 == N % 4)
!*\onslide<12->*!           System.out.print("*");
!*\onslide<13->*!       else System.out.print(" ");
!*\onslide<10->*!    }
!*\onslide<14->*!    System.out.println();
!*\onslide<9->*!}
\end{plainjava}
\endcolumns
\onslide<15->{\begin{tikzpicture}[overlay,remember picture]
    \node[above left=.55cm,xshift=.25cm] at(current page.south east) {\only<15->{\textattachfile{\curpath NestedRhombusElegant.java}{NestedRhombusElegant.java}}};
\end{tikzpicture}}
    \end{itemize}
\end{frame}

\iffull
\begin{frame}[fragile]{Can we do it in one line?}
    \begin{itemize}[<+(1)->]
        \item \textbf{Nicht klausurrelevant}, aber durch Anfrage (\only<2->{\textattachfile{\curpath NestedRhombusOneLine.java}{NestedRhombusOneLine.java}}):
\lstfs{8}
\begin{plainjava}
!*\onslide<3->*!System.out.println(args.length != 1 ? "invalid argument" :
!*\onslide<4->*!    IntStream.rangeClosed(-Integer.parseInt(args[0]) + 1, Integer.parseInt(args[0]) - 1)
!*\onslide<5->*!    .boxed()
!*\onslide<6->*!    .map(s -> Map.entry(s, !*\onslide<7->*!IntStream.rangeClosed(
!*\onslide<7->*!           -Integer.parseInt(args[0]) + 1,
!*\onslide<7->*!           Integer.parseInt(args[0]) - 1).boxed()!*\onslide<6->*!))
!*\onslide<7->*!    .map(s -> s.getValue()
!*\onslide<8->*!        .map(x -> !*\onslide<9->*!Math.abs(s.getKey()) + Math.abs(x) <= (Integer.parseInt(args[0]) - 1) &&
!*\onslide<9->*!            (Math.abs(s.getKey()) + Math.abs(x)) % 4 == (Integer.parseInt(args[0]) - 1) % 4
!*\onslide<9->*!            ? "*" : " "
!*\onslide<8->*!        )!*\onslide<10->*!.reduce("", (a, b) -> a + b)
!*\onslide<7->*!    )
!*\onslide<11->*!    .collect(Collectors.joining("\n"))
!*\onslide<3->*!);
\end{plainjava}
    \end{itemize}
\end{frame}


\begin{frame}[fragile]{Übungsblatt 5 - Zusatzaufgabe\hfill Weiter Kommentare}
    \begin{itemize}[<+(1)->]
        \itemsep6pt
        \item Man kann das Problem auch leicht rekursiv lösen. Das war so aber nicht gedacht.\pause \infoblock{Man kann jede Zeile auf die Basisfälle \say{Nichts}, \say{\bjava{*}}, \say{\bjava{* *}},~\ldots\ reduzieren und somit rekursiv bauen. Auch ne nette one-liner Idee\ldots\ Was sind aber die Nachteile der Rekursion?}
        \item Was würde sich ändern, wenn wir andere Formen wie eine Ellipse möchten oder gar mehrere Formen kombinieren?
        \item Ich kann Zusatzaufgaben wie diese nur wärmstens empfehlen. Versucht es!
        \item Wer sich unterfordert fühlt, kann immer noch nach Handicaps oder weiteren Aufgaben fragen.
        \item Code-Kommentare sollten nicht redundanterweise erklären, was die jeweilige Zeile macht, sondern maximal was die Idee/das Konzept dahinter war. Kommentare sind an sich auch Code! Nur für menschliche \say{Compiler}. \info{Oder sowas wie GPT-3.}
    \end{itemize}
\end{frame}
\fi

\iffull
\section{Altes Gedönse}

\subsection{Alte Präsenzaufgabe -- Was passiert?}
\begin{frame}[fragile, c]{Präsenzaufgabe, zum Auflockern}
    \begin{aufgabe}{Was tun Sie?!}
~\\[-\baselineskip]\begin{minipage}{0.475\linewidth}
    \pause{}Was produziert der rechts abgebildete Code, wobei:
        \pause{}{\lstfs{9}\begin{plainjava}
static int a = 5;
static int b = 10;
static boolean c = true;
        \end{plainjava}}
\end{minipage}\quad\pause{}\qquad\begin{minipage}{0.5\linewidth}
    ~\lstfs{7}\begin{plainjava}
public static void main(String[] args) {
    methode1(b, a);
    methode2(a, c);
}

static void methode1(int a, int b) {
    boolean c = false;
    b = b * a;
    System.out.println(a + " " + b + " " + c);
}

static void methode2(int b, boolean d) {
    b = a - b;
    System.out.println(a + " " + b + " " +
                        c + " " + d);
}
\end{plainjava}
\end{minipage}
\end{aufgabe}
\end{frame}

\begin{frame}[t,fragile]{Präsenzaufgabe - Lösung}
    \begin{itemize}[<+(1)->]
        \item Die Ausgabe lautet:\pause{}
\begin{java}[style=lecture@BashOutput]
10 50 false
5 0 true true
\end{java}
        \item Warum?\pause{} \bjava{false} kommt über das lokale \T{c},\pause{} die \bjava{10} über das lokale \T{a}, welches als Parameter den Wert des globalen \T{b} erhält.\pause{} Die \bjava{50} über \bjava{b * a}, wobei hier die globalen Werte verwendet werden\pause{} \info{allerdings heißt \T{a} jetzt \T{b} und \T{b} jetzt \T{a}}.
        \item Warum?\pause{}  Die \bjava{5} kommt durch das globale \T{a}, da es kein lokales Pendant besitzt.\pause{} Die \bjava{0} daher,\pause{} dass das \emph{lokale} \T{b} über die Parameter den Wert des globalen \T{a} trägt\pause{} und somit (informal) \bjava{b = a - a} gerechnet wird.\pause{} Das \T{d} entspricht dem Wert des globalen \T{c},\pause{} welches in \T{methode1} nicht überschrieben sonder überlagert wurde.
    \end{itemize}
\end{frame}

\subsection{Präzisierung: Überlagern}
\begin{frame}[t,fragile]{Präzise Antworten: Überlagerung}
    \begin{itemize}[<+(1)->]
        \item In Java gibt es verschiedene Sichtbarkeiten. \pause{} Wir werden uns hier ansehen, was bei gleichnamigen Variablen geschieht.
        \item Wir werden in diesem Kontext \emph{nicht} auf den Stack eingehen.
        \item Grundlegend sprechen wir von \emph{globaler} und \emph{lokaler} Sichtbarkeit einer Variable:\pause{}{\lstfs{9}
\begin{plainjava}
// Global:
!*\pause*!public static int globaleZahl; // !*\solGet{comments}{$\leftarrow$}*! Überall wo Klasse sichtbar
!*\pause*!private static int klasseZahl; // !*\solGet{comments}{$\leftarrow$}*! Überall innerhalb der Klasse
!*\pause*!public static int main(String[] args){
!*\pause*!    // Lokal:
    int lokaleZahl;!*\pause*! // !*\solGet{comments}{$\leftarrow$}*! Nur innerhalb von main
}
\end{plainjava}}
    \item Nun mal sehen, wie die Variablen überschrieben werden.
    \end{itemize}
\end{frame}

\rExecute{javac Verdraengung.java}
\begin{frame}[t,fragile]{Verdrängungs'v'reuden}
    \begin{itemize}[<+(1)->]
        \item Es können in Methoden nur globale Variablen überschrieben werden.
        \item Überschrieben werden sie entweder durch die Erstellung einer gleichnamige Variable\pause{} oder einem Parameter mit gleichem Namen.
        \item Der Typ spielt \emph{keine} Rolle. So erzeugt:\pause{}
        {\lstfs{11}%
        \ijava[firstline=3,lastline=7,firstnumber=3]{Verdraengung.java}
        }
        \pause{}die Ausgabe \rbash{java Verdraengung}
    \end{itemize}
\end{frame}

\begin{frame}[t,fragile]{Super-Lokal mit Blöcken}
    \begin{itemize}[<+(1)->]
        \item Wir können in Java (fast) überall geschwungene Klammerpaare setzen. \pause{}
              Sie fassen Anweisungen wieder zu einem lokalen Block zusammen.
        \item In diesem Block existieren alle Variablen der äußeren Blöcke,\pause{} wir können also keine neuen erstellen. Allerdings überleben die Variablen nur den Block.
        \item Zuerstmal ein eher sinnfreies Beispiel\ldots{}
    \end{itemize}
\end{frame}

\rExecute{javac Block.java}
\begin{frame}[t,fragile]{Super-Lokal mit Blöcken}
    \begin{itemize}[<+(1)->]
        \item Der Code
{\lstfs{8}\begin{plainjava}
!*\pause*!static int test = 0;
!*\pause*!public static void main(String[] args) {
    !*\pause*!System.out.print(test + " ");
    !*\pause*!{
        int test = 15; // Überlagere lokal
        System.out.print(test + " ");
    }
    !*\pause*!System.out.print(test);
}
\end{plainjava}}
        \pause{}erzeugt die Ausgabe \rbash[:\pause{}]{java Block}
        \item Wann ist dies von Interesse?\pause{} Bei \bjava{if} zum Beispiel \ldots
    \end{itemize}
\end{frame}

\begin{frame}[t,fragile]{Wie, wir kannten Blöcke schon?}
    \begin{itemize}[<+(1)->]
        \item Wenn wir nach einem \bjava{if}/\bjava{while}/\ldots{} Kopf eine geschwungene Klammer setzen,\pause{} starten wir einen neuen Block.
        \item Hinweis: Bei Methoden, Klassen, Enums, \ldots{} gehören die Klammern dazu (und sind deswegen anders).
        \item Allgmein noch: \pause{}\textit{achtet auf eine präzise Beschreibung!}\pause{} \say{Java weiß dann nicht was es tun soll} klingt zwar süß,\pause{} aber erklärt nicht warum.
    \end{itemize}
\end{frame}

\subsection{Präsenzaufgabe -- Vektoren}
\begin{frame}{Präsenzaufgabe, nun Vektoriell}
    \begin{itemize}[<+(1)->]
        \item Im Folgenden soll eine Vektor Klasse konstruiert werden.
        \item Zu unterstützen sind \(n\) Dimensionen \(\R^n\) (\(n \in \N\), \(n > 0\)).
        \item Ein Objekt der Vektor-Klasse hat die folgenden Informationen zu halten:
\begin{description}[Vektorkomponenten:]
    \item[Vektorkomponenten:] Eine Folge von \(n\) Werten, die nach der Instanziierung alle \emph{nicht mehr} verändert werden dürfen.
    \item[Betrag des Vektors:] Der Betrag des Vektors (ebenfalls unveränderbar).
\end{description}
        \item Um die Präsenzaufgabe übersichtlicher zu gestalten, teile ich sie in mehrere auf.
    \end{itemize}
\end{frame}

\begin{frame}[fragile, c]{Vectorias, the great one}
\begin{aufgabe}{We please the mighty}
Für einen unveränderbaren Vektor \(\R^n\) (\(n \in \N\), \(n > 0\)): \begin{enumerate}[<+(1)->]
    \item Schreiben Sie eine Klasse \bjava{Vector}, welche alle notwendigen Attribute (unveränderbare Vektorkomponenten, unveränderbarer Betrag), enthält.
    \item Schreiben Sie eine \bjava{public String toString()}-Methode. Die nach Java-Standard alle Attribute ausgibt.\pause{} Hier eine beispielhafte Ausgabe. Das genau Format ist Ihnen überlassen:
\begin{plainjava}
{ magnitude: 2.0, values: [0.0, 2.0, 0.0, 0.0, 0.0] }
\end{plainjava}
\end{enumerate}
\onslide<1->
\end{aufgabe}
\end{frame}

\begin{frame}[fragile, c]{Vectorias, the great second}
\begin{aufgabe}{We cheese the mighty}
Für einen unveränderbaren Vektor \(\R^n\) (\(n \in \N\), \(n > 0\)): \begin{enumerate}[<+(1)->]
    \setcounter{enumi}{2}
    \item Implementieren Sie \bjava{public static Vector create(double... values)}, welche die Einschränkungen prüft. Sind Sie erfüllt, soll ein neuer unveränderbarer Vektor und sonst \bjava{null} zurückgegeben werden.\pause{} Beispiele:
\begin{plainjava}
Vector a = Vector.create(0, 2, 0, 0, 0);
Vector b = Vector.create(13, 2.5, 3.2);
System.out.println(a); // :yields: Ausgabe von toString()
\end{plainjava}
\end{enumerate}
\onslide<1->
\end{aufgabe}
\end{frame}

\begin{frame}[fragile, c]{Vectorias, the great third}
\begin{aufgabe}{We freeze the mighty}
\pause Für einen unveränderbaren Vektor \(\R^n\) (\(n \in \N\), \(n > 0\)): \begin{enumerate}[<+(1)->]
    \setcounter{enumi}{3}
    \item Implementieren Sie eine Methode \bjava{public double[] getValues()}, welche die übergebenen Werte zurückliefert, aber die Unveränderbarkeit sicher stellt.\pause{}
\begin{plainjava}
Vector a = Vector.create(0, 2, 0);
double[] v = a.getValues();
v[0] = 3;
System.out.println(Arrays.toString(v)); // :yields: [3.0, 2.0, 0.0]
System.out.println(a); // :yields: original vector
\end{plainjava}
\end{enumerate}
\onslide<1->
\end{aufgabe}
\end{frame}


\begin{frame}[fragile, c]{Vectorias, the great fourth}
\begin{aufgabe}{We freeze the brighty}
Für einen unveränderbaren Vektor \(\R^n\) (\(n \in \N\), \(n > 0\)): \begin{enumerate}[<+(1)->]
    \setcounter{enumi}{4}
    \item Implementieren Sie eine Methode \bjava{public Vector normalize()}, welche einen neuen normalisierten Vektor erzeugt.\pause{}
    Dieser hat die selbe Richtung, wie das Original, aber einen Betrag von \(1\). Die Normalisierung kann durch die folgende Formel erfolgen: \begin{equation*}
        \vec{n} = \frac{\vec{v}}{||\vec{v}||}
    \end{equation*}
\end{enumerate}
\onslide<1->
\end{aufgabe}
\end{frame}

\begin{frame}[fragile]{Vectori-bras, the great solver}
    \begin{itemize}[<+(1)->]
        \item Hier die Klasse an sich: \only<2->{\textattachfile{\curpath Vector.java}{Vector.java}}
        \item Zunächst die unveränderlichen Attribute:
\begin{plainjava}
!*\onslide<4->*!public class Vector {
!*\onslide<5->*!    private final double[] values; // darf nicht public sein
!*\onslide<5->*!    public final double magnitude; // oder mit getter
!*\onslide<4->*!}
\end{plainjava}
        \item<6-> Die \bjava{toString()}-Methode:
\begin{plainjava}
!*\onslide<7->*!public String toString() {
!*\onslide<8->*!    return "{ magnitude: " + magnitude
!*\onslide<8->*!         + ", values: " + Arrays.toString(values)
!*\onslide<8->*!         + " }";
!*\onslide<7->*!}
\end{plainjava}
    \end{itemize}
\end{frame}

\begin{frame}[fragile]{Vectori-bras, the late solver}
    \begin{itemize}[<+(1)->]
        \item Die \bjava{create(double...)}-Methode:
\begin{plainjava}
!*\onslide<6->*!// expects: valid vector (auch möglich: Exceptions)
!*\onslide<6->*!private Vector(double[] values) {
!*\onslide<7->*!    this.values = values; // to be completely sure: another copy
!*\onslide<7->*!    this.magnitude = calculateMagnitude(values);
!*\onslide<6->*!}
!*\onslide<3->*!
!*\onslide<3->*!public static Vector create(double... values) {
!*\onslide<4->*!    if(values.length <= 0)
!*\onslide<4->*!        return null;
!*\onslide<5->*!    return new Vector(values);
!*\onslide<3->*!}
\end{plainjava}
    \end{itemize}
\end{frame}


\begin{frame}[fragile]{Vectori-bras, the bait resolver}
    \begin{itemize}[<+(1)->]
        \item Hier noch die Hilfsmethode für den Betrag:
\begin{plainjava}
!*\onslide<3->*!private static double calculateMagnitude(double[] values) {
!*\onslide<4->*!    double sumOfSquares = 0;
!*\onslide<5->*!    for (double value : values) {
!*\onslide<6->*!        sumOfSquares += Math.pow(value, 2);
!*\onslide<5->*!    }
!*\onslide<7->*!    return Math.sqrt(sumOfSquares);
!*\onslide<3->*!}
\end{plainjava}
    \end{itemize}
\end{frame}

\begin{frame}[fragile]{Vectori-was? The quest revolver}
    \begin{itemize}[<+(1)->]
        \item Zum Erhalt der Werte erzeugen wir eine Kopie:
\begin{plainjava}
!*\onslide<3->*!public double[] getValues() {
!*\onslide<4->*!    return copyArray(this.values);
!*\onslide<3->*!}
!*\onslide<3->*!
!*\onslide<5->*!private double[] copyArray(double[] arr) {
!*\onslide<6->*!    double[] copy = new double[arr.length];
!*\onslide<7->*!    for (int i = 0; i < copy.length; i++) {
!*\onslide<8->*!        copy[i] = arr[i];
!*\onslide<7->*!    }
!*\onslide<9->*!    return copy;
!*\onslide<5->*!}
\end{plainjava}
    \end{itemize}
\end{frame}

\begin{frame}[fragile]{Vectori-das, the final wolf}
    \begin{itemize}[<+(1)->]
        \item Nun noch die Normalisierung:
\begin{plainjava}
!*\onslide<3->*!public Vector normalize() {
!*\onslide<4->*!    double[] normalizedValues = new double[this.values.length];
!*\onslide<5->*!    for (int i = 0; i < normalizedValues.length; i++) {
!*\onslide<6->*!        normalizedValues[i] = this.values[i] / this.magnitude;
!*\onslide<5->*!    }
!*\onslide<7->*!    return Vector.create(normalizedValues);
!*\onslide<3->*!}
\end{plainjava}
    \end{itemize}
\end{frame}


\subsection{Alte Zusatzaufgabe}
\begin{frame}[fragile]{Alte Zusatzaufgabe}
    \begin{itemize}[<+(1)->]
        \item Schach! \begin{center}
\only<2|handout:0>{%
\resizebox{!}{5.5cm}{\begin{em-chess*}%
\emChessSetBoard{%
    {black:rookB, black:knightB, black:bishopB, black:queen, black:king, black:bishopA, black:knightA, black:rookA},
    {black,black,black,black,black,black,black,black},
    {},
    {},
    {},
    {},
    {white,white,white,white,white,white,white,white},
    {white:rookA, white:knightA, white:bishopA, white:queen, white:king, white:bishopB, white:knightB, white:rookB}%
}%
\end{em-chess*}%
}}\only<3->{%
\colorlet{chessblack}{paletteA}
\resizebox{!}{5.5cm}{\begin{em-chess*}%
\emChessSetBoard{%
    {,,black:rookA},
    {,,,,,black:rookA},
    {,,,,,,,black:rookA},
    {black:rookA},
    {,,,black:rookA},
    {,,,,,,black:rookA},
    {,,,,black:rookA},
    {,black:rookA}%
}%
\end{em-chess*}%
}}%
        \end{center}
    \end{itemize}
\end{frame}

\begin{frame}[fragile]{Alte Zusatzaufgabe}
    \begin{itemize}[<+(1)->]
        \item Für Türme lässt sich die Aufgabe leicht umformulieren (\only<2->{\textattachfile{\curpath Zusatz.java}{Zusatz.java}}):\par
              \say{In jeder Zeile und jeder Spalte darf es maximal einen Turm geben}
        \item Das Grundgerüst:
\begin{plainjava}
!*\pause*!boolean[][] field = { /* :ldots: */ }
!*\pause*!boolean validField = true;
!*\pause*!for (int i = 0; i < field.length; i++) {
!*\pause*!    // :ldots:
}
\end{plainjava}
        \item Kein \say{line}, da wir damit Zeilen und Spalten abgehen werden\ldots
        \item Hinweis: Wenn wir auch prüfen möchten, ob es wirklich \(n\) Türme sind, dann können wir die Aussage ersetzen durch \say{\ldots\ genau einen Turm},\pause{} und die Türme beispielsweise zählen.
    \end{itemize}
\end{frame}

\begin{frame}[fragile,c]{Alte Zusatzaufgabe}
\begin{plainjava}
!*\pause*!// für jedes i
!*\pause*!int rooksHorizontal = 0;
!*\pause*!int rooksVertical = 0;
!*\pause*!for (int j = 0; j < field.length; j++) { // da  w = h
!*\pause*!    if(field[i][j])
!*\pause*!        rooksHorizontal += 1;
!*\pause*!    if(field[j][i])
!*\pause*!        rooksVertical += 1;
!*\pause*!}
!*\pause*!if(rooksHorizontal > 1 || rooksVertical > 1) {
!*\pause*!    validField = false;
!*\pause*!    break;
!*\pause*!}
\end{plainjava}
\end{frame}

\begin{frame}[fragile]{Alte Zusatzaufgabe}
    \begin{itemize}[<+(1)->]
        \item Zuletzt noch die Auswertung:
\begin{plainjava}
!*\pause*!if(validField)
!*\pause*!    System.out.println("Feld korrekt");
!*\pause*!else
!*\pause*!    System.out.println("Feld inkorrekt");
\end{plainjava}
    \end{itemize}
\end{frame}

\begin{frame}[fragile]{Alte Zusatzaufgabe, Alternative}
    \begin{itemize}[<+(1)->]
        \item Betrachten wir eine Prüfroutine für ein gegebenes Feld (\only<2->{\textattachfile{\curpath ZusatzB.java}{ZusatzB.java}})
\begin{plainjava}
!*\onslide<3->*!boolean isRookPositionValid(int x, int y, boolean[][] field) {
!*\onslide<4->*!    for (int i = 0; i < field.length; i++) {
!*\onslide<5->*!        if(i != x && field[y][i]) return false;
!*\onslide<6->*!        if(i != y && field[i][x]) return false;
!*\onslide<4->*!    }
!*\onslide<7->*!    return true;
!*\onslide<3->*!}
\end{plainjava}
    \end{itemize}
\end{frame}

\begin{frame}[fragile]{Alte Zusatzaufgabe, Alternative}
    \begin{itemize}[<+(1)->]
        \item Nun können wir über das Feld iterieren und alle Turmpositionen prüfen:
\begin{plainjava}
!*\onslide<3->*!public static boolean checkField(boolean[][] field) {
!*\onslide<4->*!    for(int y = 0; y < field.length; y++) {
!*\onslide<5->*!        for(int x = 0; x < field.length; x++) {
!*\onslide<6->*!            if(field[y][x] && !isRookPositionValid(x, y, field))
!*\onslide<6->*!                return false;
!*\onslide<5->*!        }
!*\onslide<4->*!    }
!*\onslide<7->*!    return true;
!*\onslide<3->*!}
\end{plainjava}
        \item<8-> Auch hier lässt sich die Prüfung beispielsweise durch einen Zähler erweitern.
    \end{itemize}
\end{frame}

\section{Ein Haufen Exkurse}
\subsection{Exkurs: Objektorientierung}
\begin{frame}[t,fragile]{Ich bin eine Biene.}
    \begin{itemize}[<+(1)->]
        \item Klassen an sich schon bekannt.\pause{} Keywords: \bjava{class} so zum Beispiel bei \T{Hamster.java}:\pause{}
        \begin{java}[morekeywords={[2]{Hamster}}]
class Hamster {
    // ...
}
        \end{java}
        \item Mittels \bjava{new} lässt sich ein neues Objekt der Klasse erstellen.
        \item Beispiel: \bjava[morekeywords={[2]{Hamster}}]{Hamster günther = new Hamster();}.\pause{} Erzeugt ein \emph{Objekt} der \emph{Klasse} Hamster mit dem Namen \T{günther}.\pause{} \info{Die Deklaration der Variable verläuft analog zu \bjava{int Hamster = ...}}
    \end{itemize}
\end{frame}

\begin{frame}[t,fragile]{Objektorientierung - Attribute}
    \begin{itemize}[<+(1)->]
        \item Dieser Klasse können wir Variablen zuweisen. \pause{} Diese besitzt dann auch jedes Objekt für sich.
        \item Um auf die Variablen von außen zugreifen zu können, deklarieren wir sie mit \bjava{public}:\pause{}
\begin{java}[morekeywords={[2]{Hamster}}]
class Hamster {
    public String name;
    public int alter;
}
\end{java}
        \item Mittels \bjava{static} deklarieren wir eine Variable, die für alle Hamster identisch ist.\pause{} \info{Deswegen auch \bjava{public static void main(...)}.} \pause{}
        \begin{java*}
public static boolean flauschig = true;
        \end{java*}
    \end{itemize}
\end{frame}

\begin{frame}[t,fragile]{Objektorientierung - Methoden}
    \begin{itemize}[<+(1)->]
        \item Analog zu \emph{main} können wir Methoden deklarieren.
        \item Nebst den Zugriffspezifikatoren (\bjava{public}, \ldots)\pause{} und dem Rückgabetyp (\bjava{void} $\widehat{=}$ nichts)\pause{} schreiben wir noch den Namen (\bjava{main}),\pause{} so wie in Klammern, welche Argumente die Methode erhält (\bjava{String[] args}).
        \item Beispiel:\pause{}
\begin{java*}
public int add(int a, int b){ /* ... */ }
\end{java*}
        Diese Methode erhält zwei Integer und liefert einen zurück. Beispiel:\pause{}
\begin{java*}
public int add(int a, int b){
    return a + b;
}
\end{java*}
        \pause{}Der Aufruf \bjava{add(2,3)} liefert uns also $5$ zurück.
    \end{itemize}
\end{frame}

\begin{frame}[t,fragile]{Ein höflicher Hamster}
    \begin{itemize}[<+(1)->]
        \item Lassen wir den Hamster sich vorstellen:\pause{}
\begin{java}[morekeywords={[2]{Hamster}},morekeywords={[6]{vorstellen}}]
class Hamster {
    public String name;
    public int alter;
    public static boolean flauschig = true;

    public void vorstellen(){
        System.out.println("Ich bin " + name + " und " + alter + " Jahre alt");
    }
}
\end{java}
    \end{itemize}
\end{frame}

\begin{frame}[t,fragile]{Objektorientierung - Konstruktoren}
    \begin{itemize}[<+(1)->]
        \item Beim Erzeugen einer Klasse durch \bjava{new} wird der Konstruktor aufgerufen.\pause{} \info{An sich handelt es sich um eine Methode ohne Rückgabetyp}.
        \item So kann der Hamster bei der Erstellung einen Namen und ein Alter erhalten:\pause{}
\begin{java}[morekeywords={[2]{Hamster}}]
// ...
public Hamster(String _name, int _alter) {
    name = _name;
    alter = _alter;
}
\end{java}
        \item Nennen wir die Variablen gleich, so werden die des Hamsters über \bjava{this} erreicht,\pause{} also: \bjava{this.name = name;}.
    \end{itemize}
\end{frame}

\begin{frame}[t,fragile]{Die Hamster nutzen~\,\smaller\Smiley}
    \begin{itemize}[<+(1)->]
        \item Erzeugen wir einmal zwei Hamster in der main-Methode:\pause{}
\begin{java}[morekeywords={[2]{Hamster}}]
Hamster dieter = new Hamster("Dieter", 15);
Hamster jenny  = new Hamster("Jennifer", 7);
\end{java}
        \pause{}Nun können wir Sie sich vorstellen lassen:\pause
\begin{java}[morekeywords={[6]{vorstellen}}]
dieter.vorstellen();
jenny.vorstellen();
\end{java}
        \pause{}Dies liefert jeweils:\pause{}
        \begin{plainvoid}
Ich bin Dieter und 15 Jahre alt
Ich bin Jennifer und 7 Jahre alt
        \end{plainvoid}
    \end{itemize}
\end{frame}

\begin{frame}[t,fragile]{Allgemeines zu Klassen}
    \begin{itemize}[<+(1)->]
        \item Jede Klasse gehört in eine eigene \T{.java}-Datei mit entsprechendem Namen.
        \item Beim Kompilieren einfach alle \T{.java}-Dateien angeben.
        \item Beim Ausführen nur die Auswählen, die die \emph{main}-Methode enthält/starten soll.
        \item Andere Klassen binden wir mittels \bjava{import} ein.\pause{} Beispiel: \pause{}
\begin{java*}
import java.util.Scanner;
\end{java*}
        \item Ab jetzt steht \bjava{Scanner} zur Verfügung.
        \item So können wir auch unsere eigene Klassen einbinden.
    \end{itemize}
\end{frame}

\subsection{Bonus: Überladen und Co}

\begin{frame}[t,fragile]{Bonus: Rückgabewerte}
    \begin{itemize}[<+(1)->]
        \item Programme liefern einen Statuswert zurück.
        \item Interpretation ist proprietär. Einzig wichtig: \begin{description}
            \item[0:] Kein Fehler.
            \item[sonst:] Zahl kodiert den Fehler.
        \end{description}
        \item Deswegen: Im Fehlerfall \bjava{System.exit(1)},\pause{} oder \bjava{System.exit(42)},\pause{} oder \ldots
    \end{itemize}
\end{frame}

\begin{frame}[t,fragile]{Bonus: Kommunikation}
    \begin{itemize}[<+(1)->]
        \item \say{\bjava{System.}} liefert drei Streams\pause{} über die kommuniziert werden kann: \begin{itemize}
            \item \bjava{System.out}: für die normale Ausgabe.
            \item \bjava{System.err}: für Fehler.
            \item \bjava{System.in}: Lesen von Nutzereingaben.
        \end{itemize}
        \item \bjava{System.out} und \bjava{System.err} funktionieren identisch.
    \end{itemize}
\end{frame}

\begin{frame}[t,fragile]{Hilfe! Ich bin Überladen}
    \begin{itemize}[<+(1)->]
        \item Methoden,\pause{} kennen wir: \pause{}
\begin{java*}
public int add (int a, int b) { /* ... */ }
\end{java*}
        \item Diese Methode besitzt eine \emph{Signatur}:\pause{} \bjava{add(int,int)}.
            Der Rückgabetyp ist \emph{nicht} Teil!
        \item Wir können mehrere Methoden mit verschiedener Signatur erstellen.
        \item Bisher bekannt: anderer Name.
        \item Aber auch: gleicher Name, andere Argumente:\pause{}
\begin{java*}
public int add (int a, int b) { /* ... */ }
!*\pause*!public int add (int a, float b) { /* ... */ }
!*\pause*!public int add (int a) { /* ... */ }
!*\pause*!public int add () { /* ... */ }
\end{java*}
    \end{itemize}
\end{frame}

\begin{frame}[t,fragile]{Prinzip: Überladen}
    \begin{itemize}[<+(1)->]
        \item Wenn wir eine Methode mit gleichem Namen definieren,\pause{}
             aber einer verschiedenen Signatur,\pause{} so heißt dies: \emph{Überladen}.
        \item Java ruft die (seiner Meinung nach!) passende Methode auf.
    \end{itemize}
\end{frame}

\begin{frame}[t,fragile]{Beispiele zur Überladung}
    \begin{itemize}[<+->]
        \item Es sei (die volle Datei: \only<1->{\textattachfile{\curpath Ueberladen.java}{Ueberladen.java}}):\pause{}
\begin{java}
public static void print(short a)   {
    System.out.println("Short: " + a);
}
\end{java}

        \item Und analog:\pause{}
\begin{java}
public static void print(int a)   {
    System.out.println("Int: " + a);
}
\end{java}
        \item Sowie für \bjava{float} und \bjava{double}.
    \end{itemize}
\end{frame}

\rExecute{javac Ueberladen.java}
\begin{frame}[t,fragile]{Beispiele zur Überladung}
    \begin{itemize}[<+(1)->]
        \item Was ergibt:\pause{}
\begin{java*}
print(2); print('a'); print(0.25); print(0.1:float:);
print((byte) 15);
\end{java*}
        \item \rbash[:\pause{}]{java Ueberladen}
    \end{itemize}
\end{frame}

\begin{frame}[t,fragile]{Wie wählt Java?}
    \begin{itemize}[<+(1)->]
        \item Die erste Ausgabe (\bjava{print(2) :yields: Int: 2}) ist logisch.\pause{} \info{Eine Konstante ist eine Zahl, in Java ein \bjava{int}.}
        \item Java kann bei \say{\bjava{print('a') :yields: Int: 97}} den \bjava{char}\pause{} durch implizite Umwandlung zu einem \bjava{int} konvertieren.
        \item \say{\bjava{print(0.25) :yields: Dou:c:ble: 0.25}}\pause{} da Fließkommazahlen in Java standardmäßig \bjava{double} sind.
        \item \say{\bjava{print(0.1:float:) :yields: Float: 0.1}}\pause{} hier
            konvertieren wir die Zahl explizit!
        \item Bei \say{\bjava{print((byte) 15) :yields: Short: 15}}\pause{}
             findet Java kein \bjava{print(byte)},\pause{} konvertiert aber nicht zu
             \bjava{print(int)} sondern sucht das \say{passendste} \bjava{print(short)}.
    \end{itemize}
\end{frame}

\subsection{Exkurs: Matrizen}
\begin{frame}[t]{Exkurs: Matrizen in Java}
    \begin{itemize}[<+(1)->]
        \item Wir kennen:\pause{} Arrays!
        \begin{center}
            \begin{tikzpicture}
                \onslide<4>{\node at(0,0) {\bjava{int[] arr = \{3, 5, 7\};}};}
                \onslide<5->{\node (n) at(0,0) {\bjava{int|custom|[]|custom| arr = \{3, 5, 7\};}};}
                \begin{scope}[every node/.append style={scale=0.7}]
                    \onslide<6->{
                        \draw[decoration={brace,mirror},decorate] (n.166)++(1pt,0) -- ++(-13pt,0) node[midway,above=3pt] {Array};
                    }
                    \onslide<7->{
                        \draw[decoration=brace,decorate] (n.190)++(2pt,0) -- ++(-17pt,0) node[midway,below=3pt] {Integer};
                    }
                    % position marker
                    \onslide<8->{
                        \draw[Kite-] (n.324) -- ++(-0.45,-0.35) node[below] {\bjava{arr|custom|[0]|custom|}};
                    }
                    \onslide<9->{
                        \draw[Kite-] (n.344) -- ++(    0,-0.35) node[below] {\bjava{arr|custom|[1]|custom|}};
                    }
                    \onslide<10->{
                        \draw[Kite-] (n.350) -- ++( 0.45,-0.35) node[below] {\bjava{arr|custom|[2]|custom|}};
                    }
                    %length
                    \onslide<11->{
                        \draw[decoration={brace,mirror},decorate] (n.8)++(-4pt,0) -- ++(-49pt,0) node[midway,above=3pt] {\bjava{arr.length :yields: 3}};
                    }
                \end{scope}
            \end{tikzpicture}
        \end{center}
        \item<12-> Wir wollen: Matrizen!
            \begin{center}
                \onslide<13->{\(\displaystyle\begin{pmatrix}
                    1 & 2 & 3 \\
                    4 & 5 & 6
                \end{pmatrix}\)}
            \end{center}
            \onslide<14->{Doch wie in Java?}
    \end{itemize}
\end{frame}

\begin{frame}[t,fragile]{Exkurs: Matrizen in Java}
    \onslide<0->{%
        \begin{tikzpicture}[remember picture,overlay]
            \node[below left=0.25cm,scale=0.7,yshift=-2cm] at(current page.north east) {\(\displaystyle\begin{pmatrix}
                1 & 2 & 3 \\
                4 & 5 & 6
            \end{pmatrix}\)};
        \end{tikzpicture}}\vspace*{-\baselineskip}%
    \begin{itemize}[<+(1)->]
        \item Recht leicht:\pause{} Array von Arrays!
        \begin{center}
            \begin{tikzpicture}
                \onslide<4>{\node at(0,0) {\bjava{int[][] arr = \{\{1,2,3\}, \{4,5,6\}\};}};}
                \onslide<5->{\node (n) at(0,0) {\bjava{int[]|custom|[]|custom| arr = \{|custom|\{1,2,3\}|custom|, |custom|\{4,5,6\}|custom|\};}};}
                \begin{scope}[every node/.append style={scale=0.7}]
                    \onslide<6->{
                        \draw[decoration={brace,mirror},decorate] (n.169)++(-9pt,0) -- ++(-14pt,0) node[midway,above=3pt] {Array};
                    }
                    \onslide<7->{
                        \draw[decoration=brace,decorate] (n.186)++(14pt,0) -- ++(-29pt,0) node[midway,below=3pt] {Arrays};
                    }
                    % position marker
                    \onslide<8->{
                        \draw[Kite-] (n.321) -- ++(0,-0.35) node[below] {\bjava{arr|custom|[0]|custom|}};
                    }
                    \onslide<9->{
                        \draw[Kite-] (n.350)++(16pt,0) -- ++(    0,-0.35) node[below] (wa) {\bjava{arr|custom|[1]|custom|}};
                    }
                    %length
                    \onslide<10->{
                        \draw[decoration={brace,mirror},decorate] (n.7)++(18.5pt,0) -- ++(-103pt,0) node[midway,above=3pt] {\bjava{arr.length :yields: 2}};
                    }
                    \onslide<11->{
                        \draw[Kite-] (wa.east)++(-2pt,-1pt) -- ++(0.5,0) node[right,align=left] {Wieder ein\\Array};
                    }
                \end{scope}
            \end{tikzpicture}
        \end{center}
        \item<12-> So gilt: \onslide<13->{\bjava{arr[0] :yields: \{1, 2, 3\}}}
        \item<14-> Mehrdimensionale Arrays sind also \onslide<15->{\say{ein Array aus Arrays von Integern}.}
        \item<16-> Kompliziert? \onslide<17->{Das Muster lässt sich durch Übung leicht einprägen.}
        \item[\color{cprimary}\tiny\raisebox{1.45pt}{$\blacksquare$}]<18-> Vorerst: \onslide<19->{\bjava{arr.length} = Anzahl Zeilen} \onslide<20->{und} \onslide<21->{\bjava{arr[0].length} = Anzahl Spalten.}
    \end{itemize}
\end{frame}

\begin{frame}[t,fragile]{Exkurs: Matrizen multiplizieren}
    \begin{itemize}[<+(1)->]
        \item Sollte eigentlich aus Mathematikveranstaltungen bekannt sein.
        \item Seien \(A, B\) Matrizen.\pause{} Genauer: \(A \in \R^{a, k}\) und \(B \in \R^{k, b}\) mit \(a, k, b \in \N\).\pause{}
        Dann gilt für \(C = A \cdot B\), dass \(C \in \R^{a, b}\).
        \item Kurz:\pause{} \say{Spalten von \(A\) = Zeilen von \(B\)}.\pause{} \info{Java: \bjava{a[0].length == b.length}}.
        \item Formal berechnet sich jeder Eintrag \(c_{z,s}\) (Zeile \(z\), Spalte \(s\)) über:\pause \[c_{z,s} = \sum_{j = 1}^k a_{z,j} \cdot b_{j,s}.\]
        \item Nicht anschaulich? Vermutlich.
    \end{itemize}
\end{frame}

% i've written it without the padding :D
\def\fakepad#1{\raisebox{-0.75em}[0pt][0pt]{#1}}
\begin{frame}[t,fragile]{Exkurs: Matrizen multiplizieren -- anschaulich}
    \begin{itemize}[<+(1)->]
        \item Seien \(A\) und \(B\):\pause{}
        \[A = \begin{pmatrix}
            1 & 3 & 4 & -1\\
            1 & 1 & 3 & 0\\
        \end{pmatrix},\pause{}\quad B = \begin{pmatrix}
            3 & 0 \\
            1 & 2 \\
            1 & 1 \\
            1 & 2
        \end{pmatrix}\]
        \item Anzahl Spalten von \(A\) entspricht Anzahl Zeilen von \(B\)\pause{} \(\Rightarrow\) \(A \cdot B = C \in \R^{2, 2}\).
        \item Zur Anschaulichkeit schreiben wir sie um!
    \end{itemize}
\end{frame}


\begin{frame}[t,fragile]{Exkurs: Matrizen multiplizieren -- anschaulich}
    \onslide<0->{%
    \begin{tikzpicture}[remember picture,overlay]
        \node[below left=0.25cm,scale=0.7,yshift=-2cm] at(current page.north east) {\(\displaystyle A = \begin{pmatrix}
            1 & 3 & 4 & -1\\
            1 & 1 & 3 & 0\\
        \end{pmatrix},\quad B = \begin{pmatrix}
            3 & 0 \\
            1 & 2 \\
            1 & 1 \\
            1 & 2
        \end{pmatrix}\)};
    \end{tikzpicture}}\vspace*{-\baselineskip}%
    \begin{itemize}[<+(1)->]
        \item Seien \(A\) und \(B\) formatiert als:\pause{}
        \[\begin{array}{ccc}
            &  \begin{pmatrix}
              3 & 0 \\
              1 & 2 \\
              1 & 1 \\
              1 & 2
            \end{pmatrix} & \onslide<4->{\textcolor{cprimary}{B}}\\[2.4em]
            \begin{pmatrix}
                1 & 3 & 4 & -1\\
                1 & 1 & 3 & 0
            \end{pmatrix} & \onslide<5->{\color{gray}\begin{pmatrix}
              a&b\\
              c&d
            \end{pmatrix}} \\
            \onslide<4->{\fakepad{\textcolor{cprimary}{A}}}
        \end{array}\]
        \item<6-> Berechnen wir zunächst \(a \,\widehat{=}\, c_{1,1}\)
        \item[]<7-> \info{Die Multiplikation ist möglich, da die linke Matrix so \say{breit}, wie die obere \say{hoch} ist.}
    \end{itemize}
\end{frame}

\def\shift{.133}
\def\nshift{.1}
\def\buff#1{\makebox[1.22ex]{$#1$}}
\begin{frame}[t,fragile]{Exkurs: Matrizen multiplizieren -- anschaulich, a}
    \onslide<0->{%
    \begin{tikzpicture}[remember picture,overlay]
        \node[below left=0.25cm,scale=0.7,yshift=-2cm] at(current page.north east) {\(\displaystyle A = \begin{pmatrix}
            1 & 3 & 4 & -1\\
            1 & 1 & 3 & 0\\
        \end{pmatrix},\quad B = \begin{pmatrix}
            3 & 0 \\
            1 & 2 \\
            1 & 1 \\
            1 & 2
        \end{pmatrix}\)};
    \end{tikzpicture}}\vspace*{-\baselineskip}%
    \begin{itemize}[<+(1)->]
        \item<1-> \(A \cdot B\):\pause{} \begin{center}
            \begin{tikzpicture}
                \node at(0,0) {\(\displaystyle\begin{array}{ccc}
                    &  \begin{pmatrix}
                      3 & 0 \\
                      1 & 2 \\
                      1 & 1 \\
                      1 & 2
                    \end{pmatrix} & {\textcolor{cprimary}{B}}\\[2.4em]
                    \begin{pmatrix}
                        1 & 3 & 4 & -1\\
                        1 & 1 & 3 & 0
                    \end{pmatrix} & {\color{gray}\begin{pmatrix}
                      \only<1-6|handout:0>{\buff{\mathbf{a}}}\only<7->{\buff{\mathbf{9}}}&\buff{b}\\
                      \buff{c}&\buff{d}
                    \end{pmatrix}} \\
                    \fakepad{\textcolor{cprimary}{A}}
                \end{array}\)};
                \draw[dashed,thin, color=lightgray] (-2.65,-0.79) -- ++(4.65,0);
              \draw[dashed,thin, color=lightgray] (1.05,2) -- ++(0,-3.5);
              \onslide<3->{
                \draw[Kite-Kite,color=paletteA] (-2.125+\shift,-0.35) to[bend left, edge node={node[above,sloped] {\(\cdot\)}}] (0.8-\nshift,1.65);
              }
              \onslide<4->{
                \draw[Kite-Kite,color=paletteB] (-1.5+\shift,-0.35) to[bend left, edge node={node[above,sloped] {\(\cdot\)}}] (0.8-\nshift,1.15);
              }
              \onslide<5->{
                \draw[Kite-Kite,color=paletteC] (-0.99+\shift,-0.35) to[bend left, edge node={node[above,sloped] {\(\cdot\)}}] (0.8-\nshift,0.65);
              }
              \onslide<6->{
                \draw[Kite-Kite,color=paletteD] (-0.35+\shift,-0.35) to[bend left, edge node={node[above,sloped] {\(\cdot\)}}] (0.8-\nshift,0.25);
              }
              \node at(6.2,0) {\begin{minipage}{0.4\linewidth}
                Für \(a\) ergibt sich damit: \\
                \(a = \onslide<3->{\paletteA{1\cdot 3}} \onslide<4->{+ \paletteB{3\cdot 1}} \onslide<5->{+ \paletteC{4\cdot 1}} \onslide<6->{+ \paletteD{(-1)\cdot 1}}\)\\
                \onslide<7->{\(\hphantom{a} = 9\)}
              \end{minipage}};
              \end{tikzpicture}
        \end{center}
    \end{itemize}
\end{frame}

\begin{frame}[t,fragile]{Exkurs: Matrizen multiplizieren -- anschaulich, b}
    \onslide<0->{%
    \begin{tikzpicture}[remember picture,overlay]
        \node[below left=0.25cm,scale=0.7,yshift=-2cm] at(current page.north east) {\(\displaystyle A = \begin{pmatrix}
            1 & 3 & 4 & -1\\
            1 & 1 & 3 & 0\\
        \end{pmatrix},\quad B = \begin{pmatrix}
            3 & 0 \\
            1 & 2 \\
            1 & 1 \\
            1 & 2
        \end{pmatrix}\)};
    \end{tikzpicture}}\vspace*{-\baselineskip}%
    \begin{itemize}[<+(1)->]
        \item<1-> \(A \cdot B\):\pause{} \begin{center}
            \begin{tikzpicture}
                \node at(0,0) {\(\displaystyle\begin{array}{ccc}
                    &  \begin{pmatrix}
                      3 & 0 \\
                      1 & 2 \\
                      1 & 1 \\
                      1 & 2
                    \end{pmatrix} & {\textcolor{cprimary}{B}}\\[2.4em]
                    \begin{pmatrix}
                        1 & 3 & 4 & -1\\
                        1 & 1 & 3 & 0
                    \end{pmatrix} & {\color{gray}\begin{pmatrix}
                        \buff{9}&\only<1-6|handout:0>{\buff{\mathbf{b}}}\only<7->{\buff{\mathbf{8}}}\\
                        \buff{c}&\buff{d}
                    \end{pmatrix}} \\
                    \fakepad{\textcolor{cprimary}{A}}
                \end{array}\)};
                \draw[dashed,thin, color=lightgray] (-2.65,-0.79) -- ++(4.65,0);
              \draw[dashed,thin, color=lightgray] (1.05,2) -- ++(0,-3.5);
              \onslide<3->{
                \draw[Kite-Kite,color=paletteA] (-2.125+\shift,-0.35) to[bend left, edge node={node[above,sloped] {\(\cdot\)}}] (1.325-\nshift,1.65);
              }
              \onslide<4->{
                \draw[Kite-Kite,color=paletteB] (-1.5+\shift,-0.35) to[bend left, edge node={node[above,sloped] {\(\cdot\)}}] (1.325-\nshift,1.15);
              }
              \onslide<5->{
                \draw[Kite-Kite,color=paletteC] (-0.99+\shift,-0.35) to[bend left, edge node={node[above,sloped] {\(\cdot\)}}] (1.325-\nshift,0.65);
              }
              \onslide<6->{
                \draw[Kite-Kite,color=paletteD] (-0.35+\shift,-0.35) to[bend left, edge node={node[above,sloped] {\(\cdot\)}}] (1.325-\nshift,0.25);
              }
              \node at(6.2,0) {\begin{minipage}{0.4\linewidth}
                Für \(b\) ergibt sich damit: \\
                \(b = \onslide<3->{\paletteA{1\cdot 0}} \onslide<4->{+ \paletteB{3\cdot 2}} \onslide<5->{+ \paletteC{4\cdot 1}} \onslide<6->{+ \paletteD{(-1)\cdot 2}}\)\\
                \onslide<7->{\(\hphantom{b} = 8\)}
              \end{minipage}};
              \end{tikzpicture}
        \end{center}
    \end{itemize}
\end{frame}

\begin{frame}[t,fragile]{Exkurs: Matrizen multiplizieren -- anschaulich, c}
    \onslide<0->{%
    \begin{tikzpicture}[remember picture,overlay]
        \node[below left=0.25cm,scale=0.7,yshift=-2cm] at(current page.north east) {\(\displaystyle A = \begin{pmatrix}
            1 & 3 & 4 & -1\\
            1 & 1 & 3 & 0\\
        \end{pmatrix},\quad B = \begin{pmatrix}
            3 & 0 \\
            1 & 2 \\
            1 & 1 \\
            1 & 2
        \end{pmatrix}\)};
    \end{tikzpicture}}\vspace*{-\baselineskip}%
    \begin{itemize}[<+(1)->]
        \item<1-> \(A \cdot B\):\pause{} \begin{center}
            \begin{tikzpicture}
                \node at(0,0) {\(\displaystyle\begin{array}{ccc}
                    &  \begin{pmatrix}
                      3 & 0 \\
                      1 & 2 \\
                      1 & 1 \\
                      1 & 2
                    \end{pmatrix} & {\textcolor{cprimary}{B}}\\[2.4em]
                    \begin{pmatrix}
                        1 & 3 & 4 & -1\\
                        1 & 1 & 3 & 0
                    \end{pmatrix} & {\color{gray}\begin{pmatrix}
                        \buff{9}&\buff{8}\\
                      \only<1-6|handout:0>{\buff{\mathbf{c}}}\only<7->{\buff{\mathbf{7}}}&\buff{d}
                    \end{pmatrix}} \\
                    \fakepad{\textcolor{cprimary}{A}}
                \end{array}\)};
                \draw[dashed,thin, color=lightgray] (-2.65,-0.79) -- ++(4.65,0);
              \draw[dashed,thin, color=lightgray] (1.05,2) -- ++(0,-3.5);
              \onslide<3->{
                \draw[Kite-Kite,color=paletteA] (-2.125+\shift,-0.75) to[bend left, edge node={node[above,sloped] {\(\cdot\)}}] (0.8-\nshift,1.65);
              }
              \onslide<4->{
                \draw[Kite-Kite,color=paletteB] (-1.5+\shift,-0.75) to[bend left, edge node={node[above,sloped] {\(\cdot\)}}] (0.8-\nshift,1.15);
              }
              \onslide<5->{
                \draw[Kite-Kite,color=paletteC] (-0.99+\shift,-0.75) to[bend left, edge node={node[above,sloped] {\(\cdot\)}}] (0.8-\nshift,0.65);
              }
              \onslide<6->{
                \draw[Kite-Kite,color=paletteD] (-0.35+\shift,-0.75) to[bend left, edge node={node[above,sloped] {\(\cdot\)}}] (0.8-\nshift,0.25);
              }
              \node at(6.2,0) {\begin{minipage}{0.4\linewidth}
                Für \(c\) ergibt sich damit: \\
                \(c = \onslide<3->{\paletteA{1\cdot 3}} \onslide<4->{+ \paletteB{1\cdot 1}} \onslide<5->{+ \paletteC{3\cdot 1}} \onslide<6->{+ \paletteD{0\cdot 1}}\)\\
                \onslide<7->{\(\hphantom{c} = 7\)}
              \end{minipage}};
              \end{tikzpicture}
        \end{center}
    \end{itemize}
\end{frame}

\begin{frame}[t,fragile]{Exkurs: Matrizen multiplizieren -- anschaulich, d}
    \onslide<0->{%
    \begin{tikzpicture}[remember picture,overlay]
        \node[below left=0.25cm,scale=0.7,yshift=-2cm] at(current page.north east) {\(\displaystyle A = \begin{pmatrix}
            1 & 3 & 4 & -1\\
            1 & 1 & 3 & 0\\
        \end{pmatrix},\quad B = \begin{pmatrix}
            3 & 0 \\
            1 & 2 \\
            1 & 1 \\
            1 & 2
        \end{pmatrix}\)};
    \end{tikzpicture}}\vspace*{-\baselineskip}%
    \begin{itemize}[<+(1)->]
        \item<1-> \(A \cdot B\):\pause{} \begin{center}
            \begin{tikzpicture}
                \node at(0,0) {\(\displaystyle\begin{array}{ccc}
                    &  \begin{pmatrix}
                      3 & 0 \\
                      1 & 2 \\
                      1 & 1 \\
                      1 & 2
                    \end{pmatrix} & {\textcolor{cprimary}{B}}\\[2.4em]
                    \begin{pmatrix}
                        1 & 3 & 4 & -1\\
                        1 & 1 & 3 & 0
                    \end{pmatrix} & {\only<1-7>{\color{gray}}\begin{pmatrix}
                      \buff{9}&\buff{8}\\
                      \buff{7}&\only<1-6|handout:0>{\buff{\mathbf{d}}}\only<7>{\buff{\mathbf{5}}}\only<8-|handout:0>{\buff{5}}
                    \end{pmatrix}} \\
                    \fakepad{\textcolor{cprimary}{A}}
                \end{array}\)};
                \onslide<1-7>{\draw[dashed,thin, color=lightgray] (-2.65,-0.79) -- ++(4.65,0);}
                \onslide<1-7>{\draw[dashed,thin, color=lightgray] (1.05,2) -- ++(0,-3.5);}
              \onslide<3-7>{
                \draw[Kite-Kite,color=paletteA] (-2.125+\shift,-0.75) to[bend left, edge node={node[above,sloped] {\(\cdot\)}}] (1.325-\nshift,1.65);
              }
              \onslide<4-7>{
                \draw[Kite-Kite,color=paletteB] (-1.5+\shift,-0.75) to[bend left, edge node={node[above,sloped] {\(\cdot\)}}] (1.325-\nshift8,1.15);
              }
              \onslide<5-7>{
                \draw[Kite-Kite,color=paletteC] (-0.99+\shift,-0.75) to[bend left, edge node={node[above,sloped] {\(\cdot\)}}] (1.325-\nshift,0.65);
              }
              \onslide<6-7>{
                \draw[Kite-Kite,color=paletteD] (-0.35+\shift,-0.75) to[bend left, edge node={node[above,sloped] {\(\cdot\)}}] (1.325-\nshift,0.25);
              }
              \node at(6.2,0) {\begin{minipage}{0.4\linewidth}
                Für \(d\) ergibt sich damit: \\
                \(d = \onslide<3->{\paletteA{1\cdot 0}} \onslide<4->{+ \paletteB{1\cdot 2}} \onslide<5->{+ \paletteC{3\cdot 1}} \onslide<6->{+ \paletteD{0\cdot 2}}\)\\
                \onslide<7->{\(\hphantom{d} = 5\)}
              \end{minipage}};
              \end{tikzpicture}
        \end{center}
        \medskip
        \item<8-> Damit gilt: \(C = \begin{pmatrix}
            9 & 8 \\
            7 & 5
        \end{pmatrix}\).
    \end{itemize}
\end{frame}
\fi
% vielleicht noch was zu null
\makeatletter
\section{End-Pinguuu}
\begin{frame}[c]
\hbox{}\hfill
\onslide<2->{\begin{tikzpicture}[scale=1.15]
	\pingu[crown 2d=pingu@bronze,
	       medal=pingu@purple, tie,
	       eye patch left,
	       eye patch right,
	       right wing wave, sunglasses,
	       glow thick=pingu@yellow]
\end{tikzpicture}}\hfill\null
\end{frame}
\iffull\end{document}\fi
